\subsection{\bt{Sharing across a population of agents}{在一群智能体之间共享}}
\label{sec:population-eval}
\begin{table}[t]
\caption{\bt{Steady-state per-question cost and accuracy on TPC-DS (agents 9--16, means over \PpRuns\ runs). LLM: LLM calls of the agent; tools: tool calls; tokens include extraction.}{TPC-DS 上稳定阶段的每题代价与正确率（第 9--16 个智能体，\PpRuns\ 次运行的均值）。大模型：智能体的大模型调用次数；工具：工具调用次数；token 包括提取。}}
\label{tab:popcost}
\footnotesize
\setlength{\tabcolsep}{3pt}
% Stack bilingual labels so the numeric columns leave room for method names.
\ifbilingual\let\bt\bh\fi
\begin{tabularx}{\linewidth}{@{}Lrrrrr@{}}
\toprule
& \bt{LLM}{大模型} & \bt{Tools}{工具} & \bt{Tokens}{token} & \bt{Time}{时间} & \bt{Correct}{正确}\\
\midrule
\bt{No memory}{无记忆} & \PpNoneLateSteps & \PpNoneLateTools & \PpNoneLateTokK\,k & \PpNoneLateS\,s & \PpNoneLateAcc\%\\
\bt{Single-agent memory}{单智能体记忆} & \PpPrivateLateSteps & \PpPrivateLateTools & \PpPrivateLateTokK\,k & \PpPrivateLateS\,s & \PpPrivateLateAcc\%\\
\bt{\textbf{\system\ (ours)}}{\textbf{\system（本文）}} & \textbf{\PpMavraLateSteps} & \textbf{\PpMavraLateTools} & \textbf{\PpMavraLateTokK\,k} & \textbf{\PpMavraLateS\,s} & \textbf{\PpMavraLateAcc\%}\\
\midrule
\bt{\system\ vs.\ no memory}{\system\ 对无记忆} & & & \PpMavraLateTokXNone$\times$ & \PpMavraLateSpeedupNone$\times$ & \\
\bottomrule
\end{tabularx}
\end{table}

\bi{\emph{Protocol (Q1).} \PpAgents\ agents arrive sequentially at the TPC-DS database. Each draws seven questions from the 35-question metric workload; all methods receive the same draw within a run. Memory starts empty. A correctly answered learning question triggers extraction and publication unless it already reused a definition. Memoryless agents retain no state across questions. The single-agent baseline uses the same learning, validation, and maintenance mechanisms as \system, but restricts each agent to its own definitions, table descriptions, and join paths. This memory-scope ablation retains three configurations---no memory, private \system\ memory, and shared \system\ memory---each run \PpRuns\ times: $3\times\PpAgents\times7\times\PpRuns=1{,}680$ initial sessions and $3\times\PpAgents\times4\times\PpRuns=960$ post-restatement sessions, totaling 2,640. Published methods are compared separately in Section~\ref{sec:relbench-eval}. Initial-stage results below and post-restatement results use their respective denominators. Session tokens include any extraction they trigger.}
{\emph{设置（Q1）。}\PpAgents\ 个智能体依次到达 TPC-DS 数据库，每个从指标负载的 35 道题中抽取 7 道，同一次运行中各方法抽到的题相同。记忆初始为空；答对学习题且尚未复用定义的会话触发提取与发布。无记忆智能体不跨题保留状态。单智能体基线使用与 \system\ 相同的学习、验证和维护机制，但每个智能体只能使用自己学到的定义、表描述及连接路径。这个记忆作用域消融保留三种配置——无记忆、私有 \system\ 记忆和共享 \system\ 记忆——各运行 \PpRuns\ 次：初始阶段 $3\times\PpAgents\times7\times\PpRuns=1{,}680$ 个会话，更正后 $3\times\PpAgents\times4\times\PpRuns=960$ 个会话，合计 2,640 个。已发表方法另在第~\ref{sec:relbench-eval} 节比较。下文初始阶段结果与更正后结果分别使用对应阶段的分母；会话 token 包括其触发的提取。}

\bi{\emph{Memory-scope ablation.} Over \PpAgents\ agents, shared \system\ memory uses \PpMavraTotalTokM\,M tokens, \PpMavraTokSaveNone\% fewer than no memory, and answers \PpMavraAcc\% correctly versus \PpNoneAcc\%. Early sessions perform \PpMavraLearns\ extractions per run for five metrics; \PpMavraReused\% of all questions reuse a definition. Table~\ref{tab:popcost} reports the late-session work and accuracy. This controlled comparison varies memory scope within \system; the external-method comparison is reported in Section~\ref{sec:relbench-eval} and Figure~\ref{fig:reuse-work}.}
{\emph{记忆作用域消融。}\PpAgents\ 个智能体使用共享\system\ 记忆时共花\PpMavraTotalTokM\,M token，比无记忆少\PpMavraTokSaveNone\%，正确率为\PpMavraAcc\%，无记忆为\PpNoneAcc\%。早期会话每次运行为五个指标提取\PpMavraLearns\ 次；全部问题中\PpMavraReused\%复用定义。表~\ref{tab:popcost}报告后期会话的工作量与正确率。该受控对照改变\system\ 内部的记忆作用域；外部方法比较见第~\ref{sec:relbench-eval}节与图~\ref{fig:reuse-work}。}

\bi{\emph{Why memory must be shared.} Private memory extracts \PpPrivateLearns\ times per run but reuses on only \PpPrivateReused\% of questions: individual agents seldom revisit learned metrics. It spends \PpPrivateTotalTokM\,M tokens over \PpAgents\ agents, including \PpPrivateLateTokK\,k per late question; \system\ uses \PpMavraTokSavePrivate\% fewer tokens and runs \PpMavraSpeedupPrivate$\times$ faster.}
{\emph{为什么记忆必须共享。}私有记忆每次运行提取 \PpPrivateLearns\ 次，但仅在 \PpPrivateReused\% 的问题上复用，因为单个智能体很少重问已学指标。\PpAgents\ 个智能体共花 \PpPrivateTotalTokM\,M token，后来的智能体每题仍花 \PpPrivateLateTokK\,k；\system\ 少用 \PpMavraTokSavePrivate\% 的 token，快 \PpMavraSpeedupPrivate\ 倍。}

\bi{\emph{Publication fidelity in the metric studies.} The separate controlled metric workload also checks whether extraction preserves a verified answer. Across \PubAllJudged\ verified sessions, SQL replay rejects \PubReplayRefused\ erroneous candidates and the calculation check rejects \PubChainRefused\ more. In one example, extraction drops the electronics-category filter: a verified answer of \PubExAnswer\ becomes \PubExCanon\ under the candidate SQL, a \PubExRatio-fold error. This experiment tests the publication gate's role, independently of the RelBench token comparison.}
{\emph{指标实验中的发布保真性。}独立的受控指标负载检验提取是否保留已核实的答案。在\PubAllJudged{}个答案已核实的会话中，SQL重放拒绝\PubReplayRefused{}个错误候选，算式检查另拒绝\PubChainRefused{}个。一例中，提取丢失了电子品类过滤，使已验证答案\PubExAnswer{}变为候选SQL返回的\PubExCanon，误差达\PubExRatio{}倍。该实验检验发布检查的作用，与RelBench的token比较各自回答不同问题。}
